\section{Análisis de los planes de ejecución}
\label{sec:plans}

Spark no ejecuta las operaciones en el orden en que se escriben. Cada consulta pasa por varias fases
\cite{armbrust2015}: el analyzed logical plan (qué se pide, con nombres y tipos resueltos), el optimized logical
plan (las reglas de Catalyst lo reescriben), el physical plan (cómo se ejecuta cada operador) y, cuando llega una
action, la ejecución en stages y tasks. Esta sección analiza dos planes con \code{explain("formatted")}
(Figura~\ref{fig:plans}) y luego la ejecución completa de una consulta simple (notebook 03).

\begin{figure}[htbp]
\centering
\begin{minipage}[t]{0.585\textwidth}
\begin{lstlisting}[style=plan]
AdaptiveSparkPlan (25)
+- Sort (24)
   +- Exchange (23)
      +- Project (22)
         +- Filter (21)
            +- Window (20)
               +- WindowGroupLimit (19)
                  +- Sort (18)
                     +- Exchange (17)
                        +- WindowGroupLimit (16)
                           +- Sort (15)
                              +- Project (14)
                                 +- BroadcastHashJoin Inner BuildRight (13)
                                    :- HashAggregate (5)
                                    :  +- Exchange (4)
                                    :     +- HashAggregate (3)
                                    :        +- Filter (2)
                                    :           +- Scan parquet  (1)
                                    +- BroadcastExchange (12)
                                       +- Filter (11)
                                          +- HashAggregate (10)
                                             +- Exchange (9)
                                                +- HashAggregate (8)
                                                   +- Filter (7)
                                                      +- Scan parquet  (6)
\end{lstlisting}
\subcaption{Plan 1: top de distritos por incidencia (Q2).}
\label{fig:plan1}
\end{minipage}\hfill
\begin{minipage}[t]{0.385\textwidth}
\begin{lstlisting}[style=plan]
AdaptiveSparkPlan (9)
+- Project (8)
   +- BroadcastHashJoin LeftSemi BuildRight (7)
      :- Filter (2)
      :  +- Scan parquet  (1)
      +- BroadcastExchange (6)
         +- Project (5)
            +- Filter (4)
               +- Scan parquet  (3)
\end{lstlisting}
\subcaption{Plan 2: casos en distritos con emergencia por lluvias (Q5).}
\label{fig:plan2}
\end{minipage}
\caption{Physical plans de \code{explain("formatted")}. Solo se muestra el árbol; el número identifica a cada
operador en el detalle que imprime Spark.}
\label{fig:plans}
\end{figure}

\subsection{Plan 1: top de distritos por incidencia (Q2)}

\textbf{Logical plan.} Contar casos y sumar población por (año, distrito), unirlos con un left join, quedarse con
los distritos de 10,000 habitantes o más, calcular el ranking dentro de cada año y quedarse con el top 5.

\textbf{Optimización de Catalyst.}
\begin{enumerate}
  \item \textbf{Outer → inner join.} El filtro \code{population >= 10000} elimina las filas con nulos que
        agregaría el left join, así que Catalyst lo convierte en \code{Join Inner}.
  \item \textbf{Inferred filters y predicate pushdown.} El rango de años escrito solo para los casos llega
        también a la población a través del join key, y el filtro de población se aplica antes del join
        (operador 11).
  \item \textbf{Partition pruning.} \code{PartitionFilters: [year >= 2022, year <= 2025]}: solo se leen 4 de
        las 26 carpetas de años. El dynamic partition pruning agrega, en tiempo de ejecución, los años presentes
        en la población.
  \item \textbf{Column pruning.} El scan de casos lee \code{struct<ubigeo:string>}, 1 de 14 columnas (el año
        sale del nombre de la carpeta).
  \item \textbf{Top-N antes del shuffle.} \code{WindowGroupLimit} conserva los 5 mejores distritos por año antes
        del shuffle del window (16) y después (19).
\end{enumerate}

\textbf{Physical plan.} Cada agregación se hace en dos fases: un \code{HashAggregate} parcial en cada task, un
\code{Exchange} con \code{hashpartitioning} y un \code{HashAggregate} final (operadores 3–5 y 8–10). La población
agregada ocupa menos que el umbral de 10~MB de \code{autoBroadcastJoinThreshold}, así que se copia a cada task
(\code{BroadcastHashJoin Inner BuildRight}) y los casos no pasan por un shuffle para el join
\cite{spark-tuning}. El window necesita otro \code{Exchange} por \code{year} (17) y el \code{orderBy} final, un
\code{rangepartitioning} (23): en total, cuatro shuffles y un broadcast.

La versión en Spark SQL produce el mismo optimized logical plan y el mismo physical plan. Solo difieren el
parsed y el analyzed logical plan, porque el texto SQL se convierte en un árbol con CTEs: a Catalyst no le
importa con qué API se escribió la consulta.

\subsection{Plan 2: casos en distritos con emergencia por lluvias (Q5)}

\textbf{Logical plan.} Casos de 2023–2024 \code{left\_semi} join decretos por lluvias de 2023–2024, por
(año, distrito).

\textbf{Optimización de Catalyst.} Partition pruning (\code{year IN (2023, 2024)}: 2 de 26 carpetas), filter
pushdown al Parquet reader (\code{EqualTo(emergency\_type, rain)}, que permite saltar los row groups que no
pueden coincidir), column pruning (los casos leen solo \code{ubigeo}) y dynamic partition pruning.

\textbf{Physical plan: sin shuffle.} Los decretos filtrados son pocos, así que Spark hace broadcast de ellos como
una hash table (\code{BroadcastHashJoin LeftSemi BuildRight}). Cada task revisa sus propios casos contra esa
hash table y emite cada caso como máximo una vez, sin ningún \code{Exchange}. El inner join usa la misma
estrategia, pero emite una fila por cada decreto que coincide: 1,539,044 filas en lugar de 352,213.

\subsection{Lazy evaluation, DAG, stages y tasks}

Se usa la consulta de Q1: \code{cases.groupBy("year").agg(count).orderBy("year")}. Definirla tomó milisegundos y
no lanzó ningún job: \code{groupBy}, \code{agg} y \code{orderBy} son transformations que solo amplían el logical
plan. Recién el \code{collect()}, una action, hizo que Spark planificara y ejecutara la consulta.

La action convierte el physical plan en un DAG de stages, cortado en cada \code{Exchange}, porque un stage solo
puede empezar cuando sus datos de entrada ya fueron redistribuidos (Figura~\ref{fig:dag}). Cada stage se ejecuta
como tasks, una por partition: 9 para el scan (los archivos Parquet agrupados en 9 splits para 12 cores) y 1 en
los siguientes. Con AQE, cada stage con shuffle se lanza como su propio job, así que el \code{collect()} generó 4
jobs; los stages ya calculados aparecen como \emph{skipped} y reutilizan sus shuffle files. Dentro de un stage,
los operadores encadenados se compilan en una sola función Java (whole-stage code generation).

\begin{figure}[htbp]
\centering
\begin{tikzpicture}[
    stage/.style={draw=blue!60, fill=blue!8, rounded corners=2pt, align=center, font=\footnotesize,
                  text width=30mm, minimum height=13mm},
    arrow/.style={-{Stealth[length=2mm]}, thick, black!55},
    lbl/.style={font=\scriptsize, align=center, inner sep=1.5pt}]
  \node[stage] (s1) at (0, 0)       {\textbf{Stage 1 · 9 tasks}\\Scan parquet\\\code{partial\_count}};
  \node[stage] (s2) at (6.2, 1.1)   {\textbf{Stage 2 · 1 task}\\\code{count} final\\muestreo de sort keys};
  \node[stage] (s3) at (6.2, -1.1)  {\textbf{Stage 3 · 1 task}\\\code{count} final};
  \node[stage] (s4) at (12.4, -1.1) {\textbf{Stage 4 · 1 task}\\\code{Sort}\\26 filas al driver};
  \draw[arrow] (s1.east) -- (s2.west);
  \draw[arrow] (s1.east) -- (s3.west);
  \node[lbl] at (3.15, 0) {shuffle:\\\code{hashpartitioning}\\26 filas};
  \draw[arrow, dashed] (s2) -- node[lbl, right] {límites de\\los rangos} (s3);
  \draw[arrow] (s3) -- node[lbl, above] {shuffle:\\\code{rangepartitioning}} (s4);
\end{tikzpicture}
\caption{DAG de la consulta de Q1 con AQE: 4 stages, separados por los dos \code{Exchange} (shuffles).}
\label{fig:dag}
\end{figure}

\subsection{Shuffle}

Los casos de cada año están repartidos en archivos que leen tasks distintas. Para obtener un total por año, todos
los conteos parciales del mismo año deben llegar a la misma partition: \code{Exchange hashpartitioning(year)}
envía cada fila a la partition \code{hash(year) mod N}. Cada task escribe su salida en disco, agrupada por
partition de destino (\emph{shuffle write}), y las tasks del stage siguiente recogen su parte de cada una
(\emph{shuffle read}). Es el paso más costoso, porque implica disco, serialización y, en un cluster, red. El
\code{orderBy} necesita otro shuffle (\code{rangepartitioning}) para que cada partition tenga un rango de años.
La Tabla~\ref{tab:shuffle} compara la consulta con y sin adaptive query execution (AQE).

\begin{table}[htbp]
\centering
\caption{La consulta de Q1 con y sin AQE, según las métricas de la Spark UI.}
\label{tab:shuffle}
\small
\begin{tabular}{@{}lrr@{}}
\toprule
 & Con AQE & Sin AQE \\
\midrule
Tasks del scan & 9 & 9 \\
Filas leídas & 1,065,274 & 1,065,274 \\
Filas por shuffle (bytes) & 26 (1,664) & 26 (1,664) \\
Tasks después del shuffle & 1 + 1 + 1 & 200 + 200 + 26 \\
Tiempo del \code{collect()} & 0.14 s & 0.24 s \\
\bottomrule
\end{tabular}
\end{table}

La agregación parcial reduce el shuffle a 26 filas, una por año, en lugar de 1,065,274. Sin AQE, Spark usa las
200 partitions por defecto: 200 tasks procesan 26 filas, dos veces, más 26 tasks para el sort, casi todas vacías.
Con AQE, Spark mide el shuffle y lo reduce con coalesce a 1 partition \cite{spark-tuning}. Evitar shuffles es la
principal palanca de rendimiento: el plan 1 necesita cuatro, pero ninguno para su join gracias al broadcast, y el
plan 2 no necesita ninguno.
